\documentclass{report}

%%% Работа с языком
%\usepackage{cmap}					    % поиск в PDF				    % русские буквы в формулах
\usepackage[T2A]{fontenc}			% кодировка
\usepackage[utf8]{inputenc}             % кодировка исходного текста
\usepackage[english,ukrainian]{babel}	% локализация и переносы	    
\usepackage{indentfirst}

\usepackage[a4paper, top=25mm, bottom=25mm, left=30mm, right=30mm]{geometry}

\usepackage{amsmath,amsfonts,amssymb,amsthm,mathtools} % AMS

\usepackage{makecell}% Прикольная настройка ячеек в tabular 
%http://ctan.math.utah.edu/ctan/tex-archive/macros/latex/contrib/makecell/makecell-rus.pdf

\usepackage{diagbox}
% Разделение клеток tabular по диагонали

\usepackage{multicol} % Несколько колонок
\usepackage{wrapfig} % Картинки посреди текста
\usepackage{chngcntr} % для нумерации формул
\usepackage{enumitem} % чтоб можно было делать римскую нумерацию
\usepackage{mdwlist}
\usepackage{adjustbox}
\usepackage{xcolor} % градации цветов
\usepackage{anyfontsize}
\usepackage[psdextra]{hyperref}
\hypersetup{unicode=true}
\hypersetup{
    colorlinks=true,
    linkcolor=blue,
}
\setlist[enumerate]{nosep}

%%%% Теоремы, определения и т.д. и т.п.
\theoremstyle{plain} % Это стиль по умолчанию, его можно не переопределять.
\newtheorem{theorem}{Теорема}
\newtheorem*{theorem*}{Теорема}
\newtheorem{proposition}{Твердження}
\newtheorem*{proposition*}{Твердження}
\renewcommand{\qedsymbol}{$\blacktriangle$}
 
\theoremstyle{definition} % "Определение"
\newtheorem{definition}{Означення}[section]
\newtheorem*{definition*}{Означення}
\newtheorem*{example}{Приклад}

\theoremstyle{remark}
\newtheorem*{remark}{Зауваження}
\newtheorem*{exercise}{Вправа}

\counterwithout{equation}{chapter} % нумерация формул
\counterwithin*{equation}{section}

% красивые E, D, P, K как у Ильенко
\renewcommand{\P}{\mathbb{P}}
\newcommand{\E}{\mathbb{E}}
\newcommand{\D}{\mathbb{D}}
\newcommand{\K}{\mathbb{K}}

%%%% Картинки
\usepackage{graphicx}
\usepackage{tikz}
\usepackage{pgfplots}
\usepgfplotslibrary{fillbetween}
\pgfplotsset{compat=1.17}

\usetikzlibrary{patterns}
\usetikzlibrary{shapes.geometric}
\usetikzlibrary{arrows.meta}
\usetikzlibrary{shapes.geometric}
\usetikzlibrary{positioning,arrows}
\tikzset{
	block/.style={
		draw, 
		rectangle, 
		minimum height=1cm, 
		minimum width=3cm, align=center
	}, 
	line/.style={->,>=stealth'}}
\newcommand*\circled[1]{{\scriptsize\tikz[baseline=(char.base)]
{\node[shape=circle,draw,inner sep=0.5pt] (char) {#1};}}}
\graphicspath{{pictures/}}
\DeclareGraphicsExtensions{.pdf,.png,.jpg}

%\pagestyle{headings}
\usepackage{fancyhdr}
\pagestyle{fancy}
\renewcommand{\chaptermark}[1]{\markboth{\chaptername\ \thechapter. #1.}{}}
\renewcommand{\sectionmark}[1]{\markright{#1}{}}
\fancyhf{}
\fancyhead[C]{\chaptername~\thechapter}
\fancyhead[C]{\leftmark}
\fancyfoot[C]{\thepage}

\begin{document}
\begin{titlepage}
    \thispagestyle{empty}
    \begin{center}
        \includegraphics[width = \textwidth]{pics/kpi}
        Міністерство освіти і науки України\\
        Національний технічний університет України\\
        <<Київський політехнічний інститут ім. І. Сікорського>>\\
        Інститут прикладного системного аналізу
    \end{center}
    \vspace{30mm}
    \begin{center}
        \fontsize{22}{26}\selectfont\textbf{Доповідь} \\
        з курсу <<Числові методи оптимізації>> \\
        на тему <<Симплекс-метод>>
    \end{center}
    \vspace{30mm}
    \begin{flushright}
        \textbf{Виконали} \\ 
        студенти 3 курсу \\
        Квашук Ілья \\
        Фордуй Нікіта
    \end{flushright}
    \begin{flushright}
        \textbf{Прийняла} \\
        доцент кафедри ММСА \\
        Яковлева Алла Петрівна
    \end{flushright}
    \vspace{45mm}
    \begin{center}
        \textbf{Київ 2021}
    \end{center}
\end{titlepage}

\section{Линейное программирование}
\subsection{Задача}
\begin{equation}
	f(x)=\sum_{j=1}^n c_j x_j \to \min
\end{equation}

\begin{equation}
	\begin{cases}	
		\sum_{j=1}^n a_{ij} x_j \geq b_i, &  i=1,\dots, k\\
		\sum_{j=1}^n a_{ij} x_j= b_i, & i=k+1,\dots,m
	\end{cases}
\end{equation}
\subsection{Формы}

Общая форма:

\begin{equation}
	f(x)=\sum_{j=1}^n c_j x_j \to \min
\end{equation}

\begin{equation}
\begin{cases}	
	\sum_{j=1}^n a_{ij} x_j \geq b_i, &  i=1,\dots, k\\
	\sum_{j=1}^n a_{ij} x_j= b_i, & i=k+1,\dots,m
\end{cases}
\end{equation}

Основная форма:

\begin{equation}
	f(x)=\sum_{j=1}^n c_j x_j \to \min
\end{equation}

\begin{equation}
	\begin{cases}	
		\sum_{j=1}^n a_{ij} x_j \geq b_i, i=1,\dots, m
	\end{cases}
\end{equation}

Стандартная форма:

\begin{equation}
	f(x)=\sum_{j=1}^n c_j x_j \to \min
\end{equation}

\begin{equation}
	\begin{cases}	
		\sum_{j=1}^n a_{ij} x_j \geq b_i, &  i=1,\dots, k\\
		x_j\geq0, & j=1,\dots,n
	\end{cases}
\end{equation}

Каноническая форма: 

\begin{equation}
	f(x)=\sum_{j=1}^n c_j x_j \to \min
\end{equation}
\begin{equation}
	\begin{cases}	
		\sum_{j=1}^n a_{ij} x_j= b_i, &i=1,\dots,m\\
		x_j\geq0, & j=1,\dots,n\\
		b_i\geq0, & i=1,\dots,m
	\end{cases}
\end{equation}

Переходим от общей формы к канонической: 

\begin{equation}
	-2x_1+4x_2+3x_3 \to \max
\end{equation}

\begin{equation}
	\begin{cases}	
		x_1+2x_2+x_3\leq 1\\
		2x_1-5x_2+4x_3\geq 7\\
		x_1\geq0\\
		x_2\geq0
	\end{cases}
\end{equation}
Умножим первое уравнение на (-1):
\[2x_1-4x_2+3x_3 \to \min\]
Добавим к первому условию вспомогательну неотрицательную переменную $u_1$:
\[x_1+2x_2+x_3+u_1= 1\]
Вычтем из второго условия вспомогательную неотрицательную переменную $u_2$:
\[2x_1-5x_2+4x_3-u_2= 7\]
Представим переменную $x_3$ как $x_3^1-x_3^2$, $x_3^1\geq 0 \;x_3^2\geq 0$

\[2x_1-4x_2+3x_3^1-3x_3^2 \to \min\]
\begin{equation}
	\begin{cases}	
		x_1+2x_2+x_3^1-x_3^2+u_1= 1\\
		2x_1-5x_2+4x_3^1-4x_3^2-u_2= 7\\
		x_1,x_2,x_3^1,x_3^2,u_1,u_2\geq0
	\end{cases}
\end{equation}

\end{document}
